Unes quantes gotes petites d'oli adquireixen una càrrega negativa mentre cauen a velocitat constant a través del buit entre dues plaques horitzontals separades per una distància de 2,00~cm. Entre aquestes plaques hi ha un camp elèctric uniforme, de mòdul $5,92\times 10^{4}\ \mathrm{N\,C^{-1}}$.

\begin{apartat}{1}
Dibuixeu un esquema de la situació descrita i representeu-hi les plaques esmentades, especificant el signe de cadascuna, i els camps vectorials (gravitatori i elèctric). Calculeu la diferència de potencial entre les plaques.
\end{apartat}

\begin{apartat}{1}
Dibuixeu les forces que actuen sobre una gota de massa 2,93~pg, si la gota té una càrrega tal que fa que estigui suspesa en equilibri dins del camp elèctric esmentat. Calculeu el valor d'aquesta càrrega.
\end{apartat}

\blocetiquetat{Dada:}{$g = 9,81\ \mathrm{m\,s^{-2}}$.}
